\begin{center}
\textsc{Gauss} an \textsc{Gerling}.\quad Göttingen, 7.~November 1822.
\end{center}

\par\noindent\hspace*{\parindent}\mbox{.\ .\ .\ .\ .\ .}\quad Bei sehr flachem Einfall, bei dem sich der Strahl schließlich gar nicht mehr lenken lässt, habe ich immer mit dem \textls[120]{herrlichsten} Erfolg doppelte Reflexion einsetzen lassen, indem das Licht den Heliotropen nicht unmittelbar von der Sonne, sondern über einen an geeigneter Stelle auf dem Boden aufgestellten Handspiegel erreichte.
Ich empfehle Ihnen, diesen Kunstgriff ebenfalls anzuwenden.\quad\mbox{.\ .\ .\ .\ .\ .}

\par\noindent\hspace*{\parindent}In einem bergigen Land macht das Vermessen Freude, und zwar umso mehr, je größer die Entfernungen sind, denen beim Gebrauch der Heliotrope überhaupt keine Grenzen gesetzt sind.
Mein \textls[120]{großer} Spiegel, 1 Quadratfuß, war sogar durch den dichten Rauch des Moorbrands, der mehrere Quadratmeilen bedeckte, von Lichtenberg nach Falkenberg, fast 12 Meilen, durchgedrungen. Für die winzigen Heliotropenspiegel war dieser Rauch freilich zu dicht gewesen.\quad\mbox{.\ .\ .\ .\ .\ .}
